% =====================================================================================
% Section: Temporal Graph Representation (standalone)
%
% Required in the preamble:
%   \usepackage{amsmath,amssymb,booktabs,tikz}
%   \usetikzlibrary{arrows.meta,positioning}
%
% Citation keys used (entries in docs/paper/refs.bib, to copy into the paper's .bib):
%   nist800207, jodie, tgat, tgn, euler, lanl, cloudtrailtgn, ja3
%
% Labels defined: sec:tgr, sec:tgr-nodes, sec:tgr-unroll, sec:tgr-attr, sec:tgr-state,
%   fig:tgr-chain, tab:tgr-nodes, eq:tgr-event, eq:tgr-graph, eq:tgr-unroll,
%   eq:tgr-recency, eq:tgr-memory
% =====================================================================================

\section{Temporal Graph Representation}
\label{sec:tgr}

A Zero Trust policy decision point (PDP)~\cite{nist800207} evaluates every access
request in its full context: the network the request comes from, the client software
that issues it, the device it runs on, the identity it presents and the asset it
targets. We represent the stream of requests as a continuous-time dynamic
graph~\cite{jodie,tgn}, in which each entity is a node whose state evolves over time and
each request adds a small set of timestamped edges. Existing temporal graph detectors
model either host-to-host authentication events, as Euler does on the LANL
data~\cite{euler,lanl}, or account-to-resource interactions in cloud
logs~\cite{cloudtrailtgn}. Both reduce the request context to two endpoints. Our representation keeps the five entities observed by the
PDP as distinct typed nodes. Relations such as ``this principal has never been used
from this client'' thus become edges of the graph instead of being lost inside a
feature vector.

\subsection{Access requests as a dynamic graph}

The PDP receives requests in arrival order. Request $i$ is the tuple
\begin{equation}
  e_i = \left(t_i,\; s_i,\; c_i,\; d_i,\; u_i,\; r_i,\; \mathbf{m}_i\right),
  \label{eq:tgr-event}
\end{equation}
where $t_i$ is the arrival time, $s_i$ the client network source, $c_i$ the client TLS
configuration, $d_i$ the device, $u_i$ the authenticated principal, $r_i$ the requested
resource and $\mathbf{m}_i \in \mathbb{R}^{7}$ a vector of request attributes
(Section~\ref{sec:tgr-attr}). The graph is heterogeneous. Its node set
$V = \bigcup_{\kappa \in \mathcal{T}} V_\kappa$ is partitioned by a type map
$\tau : V \to \mathcal{T}$, with
$\mathcal{T} = \{\mathrm{src}, \mathrm{cfg}, \mathrm{dev}, \mathrm{usr}, \mathrm{res}\}$.
Edges are directed, timestamped and attributed tuples $(v, w, t, \mathbf{x})$.

Each request is unrolled into a set $E_i$ of edges that share the timestamp $t_i$
(Section~\ref{sec:tgr-unroll}). The graph against which request $i$ is evaluated is the
union of the edge sets of the earlier requests admitted to the model state:
\begin{equation}
  \mathcal{G}(t_i^-) = \bigcup_{j < i,\; j \in \mathcal{A}} E_j ,
  \label{eq:tgr-graph}
\end{equation}
where $\mathcal{A}$ is the set of requests allowed by the external policy engine. Two
properties follow from this definition. First, the representation is causal: $e_i$ is
scored against $\mathcal{G}(t_i^-)$, and its edges are inserted only afterwards, and
only if the request is admitted. Second, the content of the graph is decided outside
the detector. A request denied by the policy engine never enters the behavioural
baseline against which later requests are compared.

\subsection{Typed entities}
\label{sec:tgr-nodes}

\begin{table}[t]
\centering
\caption{Node types of the access graph. The keys of the context entities (source,
configuration, device) are namespaced by type, so equal raw identifiers of different
types never share a node.}
\label{tab:tgr-nodes}
\footnotesize
\setlength{\tabcolsep}{3pt}
\begin{tabular}{@{}l l p{0.36\columnwidth} l@{}}
\toprule
Type & Key & Observed as & Static attr. \\
\midrule
Source   & \texttt{src:}\textit{ip}    & client IP address                                  & --- \\
Config   & \texttt{conf:}\textit{ja3}  & TLS client fingerprint (JA3~\cite{ja3});
                                         \texttt{conf:guest} if none                        & --- \\
Device   & \texttt{tpm:}\,/\,\texttt{ck:}\textit{id}
                                       & TPM-attested identity or persistent device cookie  & posture tier \\
User     & principal id                & authenticated principal or anonymous visitor       & --- \\
Resource & route URI                   & normalised request path                            & sensitivity \\
\bottomrule
\end{tabular}
\end{table}

Table~\ref{tab:tgr-nodes} lists the five node types. The configuration node is the
software identity of the client. Its key is the JA3 fingerprint of the TLS
ClientHello~\cite{ja3}, which changes when the tool changes, not when the network
changes. The device node is keyed by a hardware-attested identity when a TPM is
available and by a persistent device cookie otherwise.

Each node $v$ carries a static attribute vector $\mathbf{f}_v \in \mathbb{R}^{16}$,
which is sparse and depends on the type. Devices expose their posture tier (no
certificate, certificate, certificate with TPM). Resources expose a sensitivity score
derived from their classification level in the reference access-control model. The role
and clearance of the principal are not node attributes. They are claims presented with
each request, so they are carried on the access edge (Section~\ref{sec:tgr-attr}).

The node set is open. A registry maps each external key to a slot of a pre-allocated
state table: an unseen key is admitted when it first appears, and the least recently
used slot is evicted when capacity runs out. A newly admitted node has no history, so
every node also carries a \emph{hashed identity}. Its key is hashed with BLAKE2b into
one of $B = 10^5$ buckets of a learned embedding table, giving
$\mathbf{h}_v \in \mathbb{R}^{16}$. The hash is deterministic across processes and
restarts, so a node receives the same identity at training and at serving time. An
entity that never appeared during training receives an embedding at once, without
retraining and without growing a vocabulary.

\subsection{Per-request unrolling into a causal chain}
\label{sec:tgr-unroll}

\begin{figure}[t]
\centering
\begin{tikzpicture}[
  font=\sffamily\scriptsize,
  n/.style={draw, rounded corners=2pt, align=center, inner sep=3pt, minimum height=5.5mm},
  acc/.style={->, >=Stealth, semithick},
  bind/.style={->, >=Stealth, semithick, densely dashed, gray!70!black},
  lbl/.style={font=\sffamily\tiny, inner sep=1.5pt}
]
\node[n, fill=blue!8]                     (s) {source \\ \textit{ip}};
\node[n, fill=purple!10, right=9mm of s]  (c) {config \\ \textit{ja3}};
\node[n, fill=teal!10,   right=9mm of c]  (d) {device \\ \textit{tpm / cookie}};
\node[n, fill=orange!12, below=9mm of c]  (u) {user \\ \textit{principal}};
\node[n, fill=red!8,     right=9mm of u]  (r) {resource \\ \textit{route}};
\draw[bind] (s) -- (c);
\draw[bind] (c) -- (d);
\draw[bind] (c) -- (u);
\draw[bind] (d) -- (u);
\draw[acc]  (u) -- node[lbl, below] {access, $\mathbf{m}_i$} (r);
\end{tikzpicture}
\caption{Unrolling of a single request $e_i$ into five directed edges with timestamp
$t_i$. The dashed binding edges carry an all-zero attribute vector and encode only
whether, and how recently, two entities have appeared together. The access edge carries
the request attributes $\mathbf{m}_i$. The \textit{config}\,$\rightarrow$\,\textit{user}
binding links a principal to the client software that presents it.}
\label{fig:tgr-chain}
\end{figure}

Each request produces five directed edges (Fig.~\ref{fig:tgr-chain}):
\begin{equation}
  E_i = \left\{
  \begin{aligned}
    &(s_i, c_i, t_i, \mathbf{0}),\; (c_i, d_i, t_i, \mathbf{0}),\; (c_i, u_i, t_i, \mathbf{0}),\\
    &(d_i, u_i, t_i, \mathbf{0}),\; (u_i, r_i, t_i, \mathbf{m}_i)
  \end{aligned}
  \right\}.
  \label{eq:tgr-unroll}
\end{equation}
The first four are \emph{binding} edges. They record which network carried which
client, which client ran on which device, and which principal was used from which
client and from which device. Their attribute vector is zero, so what they encode is
purely structural and temporal: whether the pair has been seen before, and when. The
last edge is the \emph{access} edge and carries the request attributes. The chain follows
the order in which the context of a request is established: network, client software,
hardware, identity, asset. The edges of a request are inserted in this fixed order in
training, in offline replay and in serving.

The binding edges make visible some request patterns that a two-endpoint graph cannot
express. A principal used from a client never associated with it before appears as a
novel \textit{config}\,$\rightarrow$\,\textit{user} edge, even when the principal and
the resource are both habitual, which is the signature of credential reuse from a
different tool. A pivot with a harvested credential appears as a
\textit{device}\,$\rightarrow$\,\textit{user} binding, and possibly an access edge, that
is unusual for the entities involved. The configuration node also complements the TLS
trust bit carried in $\mathbf{m}_i$: the bit states whether the fingerprint is trusted,
while the node states \emph{which} client is in use.

The source and device keys are optional. When one of them is missing, only its own
binding edges are dropped. Without a device identity, the
\textit{config}\,$\rightarrow$\,\textit{device} and
\textit{device}\,$\rightarrow$\,\textit{user} edges disappear, while
\textit{config}\,$\rightarrow$\,\textit{user} still links the client to the principal.
Without a source address, the \textit{source}\,$\rightarrow$\,\textit{config} edge is
dropped. The configuration node is always present: when the fingerprint cannot be
resolved, the request is attached to the generic \texttt{conf:guest} node. The same
representation therefore accepts telemetry of different richness without changes to the
model.

\subsection{Edge attributes and time}
\label{sec:tgr-attr}

The access-edge attributes $\mathbf{m}_i$ contain the TLS fingerprint trust bit, the
outputs of three intrusion-detection sensor probes, the HTTP method, and the role and
clearance of the presented identity (normalised to $[0,1]$).
All of them are known to the PDP when it authorises the request. Response-side fields,
such as the status code or the response volume, are excluded because they become
available only after the decision.

Elapsed times are mapped to a vector by a learnable harmonic encoding
$\phi(\Delta t) = \cos(\boldsymbol{\omega}\,\Delta t + \mathbf{b}) \in
\mathbb{R}^{32}$~\cite{tgat,tgn}. For each edge $(v, w)$ scored at time $t$, the recency
of the pair is
\begin{equation}
  \delta_{vw}(t) =
  \begin{cases}
    \min\!\left(t - t^{\mathrm{last}}_{vw},\; \Delta_{\max}\right) & \text{if } (v,w) \text{ was seen,}\\
    \Delta_{\max} & \text{otherwise,}
  \end{cases}
  \label{eq:tgr-recency}
\end{equation}
with $\Delta_{\max} = 7$\,days, and the recency $\delta_v(t)$ of the source entity is
defined in the same way. On the access edge the source is the principal, so its request
cadence enters the model through $\delta_{u_i}(t)$, which advances with every admitted
request, rather than as an edge attribute. The cap works as a sentinel for unseen pairs.
Encoding them with the absolute clock would produce a value that grows along the stream,
so training, validation and serving would see different encodings of the same state.
Beyond one week, a silent pair is treated as never seen.

Each edge is also described by causal interaction counters. With $n_{vw}$ the number of
admitted occurrences of the pair and $n_v$ those of the source entity, the edge carries
$\left[\log(1 + n_{vw}),\; \log(1 + n_v),\; n_{vw} / (n_v + 1)\right]$. The access edge
additionally carries the same triplet for the (device, resource) pair, which records how
habitual the resource is for the device without adding a sixth temporal edge. All counters
are updated only on admitted requests, so they can be computed at serving time without
labels.

\subsection{Temporal state}
\label{sec:tgr-state}

The graph is not stored as an edge list. Its history is summarised in two bounded
per-node structures. The first is a recurrent memory
$\mathbf{s}_v \in \mathbb{R}^{256}$, following the Temporal Graph
Network~\cite{tgn}. When $v$ takes part in an admitted edge $(v, w, t, \mathbf{x})$, a
message is formed and the memory is updated by a GRU:
\begin{equation}
  \begin{aligned}
    \boldsymbol{\mu}_v &= \left[\,\mathbf{s}_v \,\|\, \mathbf{s}_w \,\|\, \mathbf{x}
                          \,\|\, \phi(t - t_v^{-})\,\right],\\
    \mathbf{s}_v &\leftarrow \mathrm{GRU}\!\left(\overline{\boldsymbol{\mu}}_v,\; \mathbf{s}_v\right),
  \end{aligned}
  \label{eq:tgr-memory}
\end{equation}
where $t_v^{-}$ is the time of the last update of $v$ and $\overline{\boldsymbol{\mu}}_v$
is the mean of the messages that $v$ receives within an update batch.

The second structure is a temporal neighbour store. For every node it holds the last
$K = 30$ incident edges, each as a (neighbour, timestamp, attributes) triple, in a
fixed-size ring buffer. Every edge is written at both endpoints, so the neighbourhood is
symmetric although the edges are directed. The store occupies
$O(|V| \cdot K \cdot d_m)$ memory, with $d_m$ the attribute size, independently of the
length of the stream, and no external graph database is needed.

When a request is scored, the temporal neighbourhood of its endpoints is expanded for up
to three hops from the store. Each node enters the encoder with the input
$\left[\mathbf{s}_v \,\|\, \mathbf{f}_v \,\|\, \mathbf{h}_v\right] \in \mathbb{R}^{288}$,
and each stored edge $e$ from a neighbour $w$ with the attribute
$\left[\phi(t_w^{-} - t_e) \,\|\, \mathbf{x}_e\right]$, where $t_e$ is the timestamp of the
edge. Both structures advance only for
admitted requests. They are therefore a bounded summary of $\mathcal{G}(t_i^-)$ in
Eq.~\eqref{eq:tgr-graph}.
